\chapter{基本相互作用}\label{ch:g11:fundamental-interactions}

在头发上摩擦气球，它便能吸起碎纸屑——
这只是几平方厘米的橡胶，却胜过整个行星的引力。
这场轻而易举的“胜利”，正是\omterm{def:g11:fundamental-interactions:ladder}{原子}能够结合在一起的原因，也是
重\omterm{def:g11:fundamental-interactions:ladder}{原子核}终将衰变的原因（见\cref{ch:g11:nucleus-radioactivity}）。
宇宙只依赖四种\omterm{def:g11:fundamental-interactions:four}{基本相互作用}。本章将逐一认识它们，比较它们的强弱，
并为每一种指定其所属的层级。

\section{从夸克到星系的物质}

\begin{definition}[结构阶梯]\label{def:g11:fundamental-interactions:ladder}
物质以层级构成，每一层都由其下一层组装而成：三个
\emph{夸克}\index{夸克}（点状，小于\qty{e-18}{m}）束缚成
一个\emph{核子}\index{核子}——质子或中子，直径约为
\qty{e-15}{m}；核子聚集成一个\emph{原子核}\index{原子核}
（\qty{e-15}{m}至\qty{e-14}{m}）；原子核加上其电子云，便是一个
\emph{原子}\index{原子}（\qty{e-10}{m}）；原子结合成
\emph{分子}\index{分子}和晶体（\qty{e-9}{m}起），
进而组成日常物体（\qty{1}{m}）、行星（\qty{e7}{m}）、
行星系统（\qty{e11}{m}）、星系（\qty{e21}{m}）和
可观测宇宙（\qty{e26}{m}）。
\end{definition}

\begin{definition}[四种基本相互作用]\label{def:g11:fundamental-interactions:four}
\emph{基本相互作用}\index{基本相互作用}是一种
不能归结为其他力的力；共有且仅有四种：
\emph{万有引力}\index{万有引力}、\emph{电磁相互作用}\index{电磁相互作用}、
\omterm{def:g11:fundamental-interactions:strong}{强相互作用}和\omterm{def:g11:fundamental-interactions:weak}{弱相互作用}。
\end{definition}

\begin{remark}\label{rem:g11:fundamental-interactions:empty}
各层之间相隔着空隙：\omterm{def:g11:fundamental-interactions:ladder}{原子}的宽度是\omterm{def:g11:fundamental-interactions:ladder}{原子核}的 $10^{5}$ 倍——
若把\omterm{def:g11:fundamental-interactions:ladder}{原子核}放大为直径 \qty{1}{cm} 的弹珠，
电子云则有一千米宽。这些稀疏的层级需要某种力将其维系在一起，
而引力几乎从不是承担此任的力。
\end{remark}

\begin{omfigure}
\begin{tikzpicture}[font=\small, x=0.24cm]
  \draw[-Stealth] (-1,0) -- (46.5,0);
  \foreach \x/\l/\d in {0/{10^{-18}}/2pt, 3/{10^{-15}}/13pt,
                        8/{10^{-10}}/2pt, 18/{10^{0}}/2pt,
                        25/{10^{7}}/2pt, 29/{10^{11}}/13pt,
                        39/{10^{21}}/2pt, 44/{10^{26}}/13pt}
    \draw (\x,0.1) -- (\x,-0.1) node[below=\d] {$\l$};
  \foreach \x/\n in {0/夸克, 3/原子核, 8/原子, 18/人, 25/地球, 
                     29/{日地距离}, 39/星系, 44/宇宙}{
    \fill[omDef] (\x,0) circle (1.6pt);
    \node[rotate=45, anchor=west, inner sep=1.5pt] at (\x,0.24) {\n};}
\end{tikzpicture}

{\small 以 10 的幂为刻度的结构阶梯：从\omterm{def:g11:fundamental-interactions:ladder}{夸克}到可观测宇宙
跨越四十四个\omterm{def:g10:orders-of-magnitude:oom}{数量级}。尺寸以
米为\omterm{def:g10:orders-of-magnitude:unit}{单位}。}
\end{omfigure}

\section{电荷}

\begin{definition}[电荷]\label{def:g11:fundamental-interactions:charge}
\emph{电荷}\index{电荷}是物质的一种属性，
\omterm{def:g11:fundamental-interactions:four}{电磁相互作用}作用于它，正如\omterm{def:g11:fundamental-interactions:four}{万有引力}作用于质量。
其\omterm{def:g10:orders-of-magnitude:unit}{单位}为\emph{库仑}\index{库仑}
（\unit{C}），有\emph{两种电性}，即正与负，二者相加时\omterm{def:g10:inertia:compensate}{相互抵消}；
同种电荷相斥，异种电荷相吸。
\end{definition}

\begin{definition}[元电荷]\label{def:g11:fundamental-interactions:elementary}
\omterm{def:g11:fundamental-interactions:charge}{电荷}是量子化的：每一个被观测到的\omterm{def:g11:fundamental-interactions:charge}{电荷}都是
\emph{元电荷}\index{元电荷}的整数倍。
$e = \qty{1.60e-19}{C}$。质子带 $+e$，电子带 $-e$，
中子的\omterm{def:g11:fundamental-interactions:charge}{电荷}为 $0$；\omterm{def:g11:fundamental-interactions:charge}{电荷} $q$ 含有 $N = \abs{q}/e$ 个元\omterm{def:g11:fundamental-interactions:charge}{电荷}。
\end{definition}

\begin{proposition}[电荷守恒]\label{prop:g11:fundamental-interactions:conservation}
孤立系统的总\omterm{def:g11:fundamental-interactions:charge}{电荷}永远不变：\omterm{def:g11:fundamental-interactions:charge}{电荷}既不能被创生，
也不能被消灭，只能发生转移——实际上是由电子，
即最轻的\omterm{def:g11:fundamental-interactions:charge}{电荷}载体，来完成这种转移。
\end{proposition}
\admitted

\begin{remark}\label{rem:g11:fundamental-interactions:conservation}
从未观测到违反的情况；其深层原因是电磁学的一种对称性，
在本科分册中给出。
\end{remark}

\begin{definition}[摩擦起电与接触起电]\label{def:g11:fundamental-interactions:charging}
摩擦两个绝缘体，会使电子从一个转移到另一个，
使二者带上电量相等、电性相反的\omterm{def:g11:fundamental-interactions:charge}{电荷}（\emph{摩擦起电}\index{摩擦起电}）；
将带电物体与中性物体接触，会使\omterm{def:g11:fundamental-interactions:charge}{电荷}在二者间共享
（\emph{接触起电}\index{接触起电}）。
其中只有电子在移动，
总\omterm{def:g11:fundamental-interactions:charge}{电荷}保持守恒。
\end{definition}

\begin{example}[计算电子个数]\label{ex:g11:fundamental-interactions:balloon}
在头发上摩擦过的气球获得\omterm{def:g11:fundamental-interactions:charge}{电荷} $q = \qty{-1.6e-8}{C}$：它
\emph{得到}了 $N = \num{1.6e-8} / \num{1.6e-19} = \num{1.0e11}$
个电子，头发则留下恰好 $+\qty{1.6e-8}{C}$ 的\omterm{def:g11:fundamental-interactions:charge}{电荷}。
\end{example}

\section{库仑定律}

\begin{theorem}[库仑定律]\label{thm:g11:fundamental-interactions:coulomb}
两个点\omterm{def:g11:fundamental-interactions:charge}{电荷} $q_1$ 和 $q_2$，相距 $d$，二者之间
沿连线相互施加大小相等、方向相反的力——同号时为斥力，
异号时为引力——其大小
由下式给出：
\[
F = k\,\frac{\abs{q_1 q_2}}{d^{2}},
\qquad k = \qty{8.99e9}{N.m^2/C^2}.
\]
\end{theorem}
\admitted

\begin{remark}\label{rem:g11:fundamental-interactions:torsion}
\omterm{def:g11:fundamental-interactions:charge}{库仑}于 1785 年通过测量一根细扭丝的扭转，建立了该定律；
为何\omterm{def:g10:orders-of-magnitude:scinot}{指数}恰为 $2$，将在大学分册中解答。
\end{remark}

\begin{omfigure}
\begin{tikzpicture}[font=\small]
  \draw[omDef, thick, fill=omDef!20] (0,0) circle (0.3) node[black] {$+$};
  \draw[omDef, thick, fill=omDef!20] (2.6,0) circle (0.3) node[black] {$+$};
  \draw[-Stealth, omThm, very thick] (-0.3,0) -- (-1.5,0)
    node[midway, above] {$\vect F_{2/1}$};
  \draw[-Stealth, omThm, very thick] (2.9,0) -- (4.1,0)
    node[midway, above] {$\vect F_{1/2}$};
\end{tikzpicture}\qquad
\begin{tikzpicture}[font=\small]
  \draw[omDef, thick, fill=omDef!20] (0,0) circle (0.3) node[black] {$+$};
  \draw[omProp, thick, fill=omProp!25] (2.6,0) circle (0.3) node[black] {$-$};
  \draw[-Stealth, omThm, very thick] (0.3,0) -- (1.1,0)
    node[midway, above] {$\vect F_{2/1}$};
  \draw[-Stealth, omThm, very thick] (2.3,0) -- (1.5,0)
    node[midway, above] {$\vect F_{1/2}$};
\end{tikzpicture}

{\small 同号\omterm{def:g11:fundamental-interactions:charge}{电荷}相斥（左），异号\omterm{def:g11:fundamental-interactions:charge}{电荷}相吸（右）；
两力大小始终相等、方向相反。}
\end{omfigure}

\begin{remark}[与引力的形式对应]\label{rem:g11:fundamental-interactions:twin}
\omterm{def:g11:fundamental-interactions:charge}{库仑}定律在符号上与\omterm{def:g11:fundamental-interactions:four}{万有引力}定律
（见\cref{ch:g10:universal-gravitation}）完全相同，只需把质量换成\omterm{def:g11:fundamental-interactions:charge}{电荷}：
\begin{center}
\begin{tabular}{lll}
源 & 质量 $m$ & \omterm{def:g11:fundamental-interactions:charge}{电荷} $q$ \\
定律 & $F = G\,\dfrac{m_1 m_2}{d^2}$ & $F = k\,\dfrac{\abs{q_1 q_2}}{d^2}$ \\[1.5ex]
常量 & $G = \qty{6.67e-11}{N.m^2/kg^2}$ & $k = \qty{8.99e9}{N.m^2/C^2}$ \\
符号 & 始终为引力 & 可为引力或斥力 \\
\end{tabular}
\end{center}
下文的两点差异是全部内容的基础：电相互作用远强于引力，
且只有电相互作用可以\omterm{def:g10:inertia:compensate}{相互抵消}。
\end{remark}

\begin{example}[两个质子，两种力同时作用]\label{ex:g11:fundamental-interactions:protons}
两个质子（$m_p = \qty{1.67e-27}{kg}$）相距 $d = \qty{1.0e-15}{m}$，
即\omterm{def:g11:fundamental-interactions:ladder}{原子核}内的近邻。电斥力为：
\[
F_E = \num{8.99e9} \times
\frac{(\num{1.60e-19})^{2}}{(\num{1.0e-15})^{2}}
\approx \qty{2.3e2}{N}
\]
——相当于一只 $23$ \omterm{def:g10:orders-of-magnitude:unit}{千克}箱子所受的\omterm{def:g10:universal-gravitation:weight}{重力}，而承载它的只是
一个质量为 $10^{-27}$ \omterm{def:g10:orders-of-magnitude:unit}{千克}的粒子。\omterm{def:g11:fundamental-interactions:four}{万有引力}为：
$F_G = \num{6.67e-11} \times (\num{1.67e-27})^{2} /
(\num{1.0e-15})^{2} \approx \qty{1.9e-34}{N}$。比值
$F_E / F_G = k e^{2} / (G m_p^{2}) \approx \num{1.2e36}$
与距离无关（$d^{2}$ 相消）：对于基本粒子，
引力弱了约 $10^{36}$ 倍，以至于在\emph{任何}距离下都无足轻重。
\end{example}

\section{强相互作用与弱相互作用}

上例留下了一个疑问：氦核中的两个质子
以数十牛的斥力相互排斥，而氦核却依然存在。

\begin{definition}[强相互作用]\label{def:g11:fundamental-interactions:strong}
\emph{强相互作用}\index{强相互作用}将\omterm{def:g11:fundamental-interactions:ladder}{夸克}
束缚成\omterm{def:g11:fundamental-interactions:ladder}{核子}，又将\omterm{def:g11:fundamental-interactions:ladder}{核子}束缚成\omterm{def:g11:fundamental-interactions:ladder}{原子核}。在 \qty{e-15}{m} 处它是
吸引力，且远强于电斥力（\omterm{def:g11:fundamental-interactions:ladder}{核子}间可达
数千牛）；超过几个 $10^{-15}$ 米后，它便消失了：
其\emph{作用范围}\index{作用范围（相互作用的）}约为
\qty{e-15}{m}。它对质子和中子一视同仁，与\omterm{def:g11:fundamental-interactions:charge}{电荷}无关。
\end{definition}

\begin{omfigure}
\begin{tikzpicture}[font=\small]
  \draw[omDef, thick, fill=omDef!20] (0,0) circle (0.34) node[black] {p};
  \draw[omDef, thick, fill=omDef!20] (3.2,0) circle (0.34) node[black] {p};
  \draw[-Stealth, omThm, thick] (-0.34,0) -- (-1.7,0)
    node[midway, above] {$F_E \approx \qty{230}{N}$};
  \draw[-Stealth, omThm, thick] (3.54,0) -- (4.9,0) node[midway, above] {$F_E$};
  \draw[-Stealth, omMeth, very thick] (0.2,0.55) -- (1.45,0.55)
    node[midway, above] {强};
  \draw[-Stealth, omMeth, very thick] (3.0,0.55) -- (1.75,0.55)
    node[midway, above] {强};
  \draw[Stealth-Stealth] (0,-0.65) -- (3.2,-0.65)
    node[midway, below] {$d = \qty{1.0e-15}{m}$};
\end{tikzpicture}

{\small \omterm{def:g11:fundamental-interactions:ladder}{原子核}内的拉锯：在 $10^{-15}$~m 处，强
吸引力（数千牛）胜过电斥力。}
\end{omfigure}

\begin{remark}[为何原子核能够存在——以及为何重核会衰变]\label{rem:g11:fundamental-interactions:nuclei}
\omterm{def:g11:fundamental-interactions:ladder}{原子核}内的两种力随距离的变化方式不同。强“胶水”作用距离短：
一个\omterm{def:g11:fundamental-interactions:ladder}{核子}只与其邻近的少数\omterm{def:g11:fundamental-interactions:ladder}{核子}结合，范围在 \qty{e-15}{m} 以内。
电斥力作用距离长：每个质子都会排斥
\emph{其他每一个}质子。随着\omterm{def:g11:fundamental-interactions:ladder}{原子核}增大，斥力的累积速度快于“胶水”的结合力：
重核需要更多中子（中子不带电，只起“胶水”作用），
而当质子数超过约 $80$ 时，即便如此也不够——
最重的\omterm{def:g11:fundamental-interactions:ladder}{原子核}只能“借来”的时间存在（见\cref{ch:g11:nucleus-radioactivity}）。
\end{remark}

\begin{definition}[弱相互作用]\label{def:g11:fundamental-interactions:weak}
\emph{弱相互作用}\index{弱相互作用}，其\omterm{def:g11:fundamental-interactions:strong}{作用范围}更短，
亦可使中子转变为质子（反之亦然）：它驱动
\cref{ch:g11:nucleus-radioactivity}中的 $\beta$ 放射性。
\end{definition}

\section{谁统治哪一层}

\begin{method}[判断尺度]\label{met:g11:fundamental-interactions:audit}
要判断大小为 $d$ 的结构受哪种相互作用支配，需检查
\omterm{def:g11:fundamental-interactions:strong}{作用范围}与抵消情况：
\begin{enumerate}
  \item $d \leq \qty{e-14}{m}$：\emph{\omterm{def:g11:fundamental-interactions:strong}{强相互作用}}起支配作用——
        它能压过电斥力，而引力则小了 $10^{36}$ 倍。
  \item 从\omterm{def:g11:fundamental-interactions:ladder}{原子}到山脉（$\qty{e-10}{m}$ 至 $\qty{e4}{m}$）：
        强力已超出\omterm{def:g11:fundamental-interactions:strong}{作用范围}；起主导作用的是\emph{\omterm{def:g11:fundamental-interactions:four}{电磁
        相互作用}}——键合、内聚、摩擦与接触。
  \item 行星及更远尺度（$d \geq \qty{e6}{m}$）：物质在极高精度下呈电中性，
        因此\omterm{def:g11:fundamental-interactions:charge}{电荷}\omterm{def:g10:inertia:compensate}{相互抵消}；质量只有一种符号，
        总是叠加——\emph{\omterm{def:g11:fundamental-interactions:four}{万有引力}}主宰天体。
\end{enumerate}
\omterm{def:g11:fundamental-interactions:weak}{弱相互作用}不占据任何层次：它不束缚任何物质，
但主宰转变过程（$\beta$ 衰变）。
\end{method}

\begin{omfigure}
\begin{tikzpicture}[font=\small, x=0.26cm]
  \draw[-Stealth] (-1,0) -- (41.5,0);
  \foreach \x/\l in {3/{10^{-15}}, 8/{10^{-10}}, 18/{10^{0}},
                     24/{10^{6}}, 39/{10^{21}}}
    \draw (\x,0.1) -- (\x,-0.1) node[below] {$\l$};
  \draw[omProp, fill=omProp!35] (3,0.3) rectangle (4,0.75);
  \node[anchor=west, omProp] at (4.4,0.52) {强相互作用};
  \draw[omDef, fill=omDef!25] (8,0.95) rectangle (22,1.4);
  \node[omDef] at (15,1.17) {电磁相互作用};
  \draw[omThm, fill=omThm!25] (24,1.6) rectangle (39,2.05);
  \node[omThm] at (31.5,1.82) {万有引力};
\end{tikzpicture}

{\small 各层由何种相互作用主导：\omterm{def:g11:fundamental-interactions:ladder}{原子核}由\omterm{def:g11:fundamental-interactions:strong}{强相互作用}，从\omterm{def:g11:fundamental-interactions:ladder}{原子}到山脉由\omterm{def:g11:fundamental-interactions:four}{电磁相互作用}，
行星及更远处由\omterm{def:g11:fundamental-interactions:four}{万有引力}。尺寸以米为\omterm{def:g10:orders-of-magnitude:unit}{单位}。}
\end{omfigure}

\section{练习}

\begin{exercise}[$\star$]\label{exo:g11:fundamental-interactions:1}
求出各尺寸的\omterm{def:g10:orders-of-magnitude:oom}{数量级}，并指出其所属层次：
质子 \qty{1.7e-15}{m}；水\omterm{def:g11:fundamental-interactions:ladder}{分子} \qty{2.8e-10}{m}；沙粒
\qty{2e-4}{m}；地球 \qty{1.3e7}{m}；银河 \qty{9.5e20}{m}。
\end{exercise}

\begin{exercise}[$\star$]\label{exo:g11:fundamental-interactions:2}
一只与头发摩擦过的气球带有\omterm{def:g11:fundamental-interactions:charge}{电荷} $q = \qty{-4.8e-9}{C}$。它是得到
还是失去了电子？共多少个？之后头发带有多少\omterm{def:g11:fundamental-interactions:charge}{电荷}，
依据哪条定律？
\end{exercise}

\begin{exercise}[$\star$]\label{exo:g11:fundamental-interactions:3}
\omterm{def:g11:fundamental-interactions:charge}{电荷} $q_1 = \qty{+2.0e-6}{C}$ 与 $q_2 = \qty{-3.0e-6}{C}$ 相距
\qty{30}{cm}。求作用力，是吸引还是排斥？
比较 $q_1$ 与 $q_2$ 所受的力。
\end{exercise}

\begin{exercise}[$\star$]\label{exo:g11:fundamental-interactions:4}
两个\omterm{def:g11:fundamental-interactions:charge}{电荷}以力 $F_0$ 相互吸引。下列情况下，此力变为多大：
距离加倍；距离除以 $3$；一个\omterm{def:g11:fundamental-interactions:charge}{电荷}加倍；
两个\omterm{def:g11:fundamental-interactions:charge}{电荷}\emph{与}距离同时加倍？
\end{exercise}

\begin{exercise}[$\star$]\label{exo:g11:fundamental-interactions:5}
指出下列各现象所对应的相互作用：(a) 月球绕地球运行；
(b) 食盐晶体的内聚；(c) 氦
\omterm{def:g11:fundamental-interactions:ladder}{原子核}的内聚；(d) 碳-14 的 $\beta$ 衰变；(e)
\cref{exo:g11:fundamental-interactions:2}中的气球粘在墙上。
\end{exercise}

\begin{exercise}[$\star\star$]\label{exo:g11:fundamental-interactions:6}
氦\omterm{def:g11:fundamental-interactions:ladder}{原子核}中两质子相距 $d = \qty{2.0e-15}{m}$，
二者之间的\omterm{def:g11:fundamental-interactions:charge}{库仑}斥力、\omterm{def:g11:fundamental-interactions:four}{万有引力}（$m_p = \qty{1.67e-27}{kg}$）及二者之比，
求之。是什么使\omterm{def:g11:fundamental-interactions:ladder}{原子核}保持完整？
\end{exercise}

\begin{exercise}[$\star\star$]\label{exo:g11:fundamental-interactions:7}
氢\omterm{def:g11:fundamental-interactions:ladder}{原子}中，电子（$m_e = \qty{9.11e-31}{kg}$）位于
距质子 $d = \qty{5.3e-11}{m}$ 处（$m_p = \qty{1.67e-27}{kg}$）。
求二者之间的\omterm{def:g11:fundamental-interactions:charge}{库仑}力与\omterm{def:g11:fundamental-interactions:four}{万有引力}，以及它们的
比值。是什么把\omterm{def:g11:fundamental-interactions:ladder}{原子}维系在一起？
\end{exercise}

\begin{exercise}[$\star\star$]\label{exo:g11:fundamental-interactions:8}
小球 A 带电 $q_A = \qty{+8.0e-9}{C}$；与之相同的小球 B
呈电中性。二者接触后，再分开至相距 \qty{10}{cm}。
两球各带多少\omterm{def:g11:fundamental-interactions:charge}{电荷}（检验\omterm{def:g11:fundamental-interactions:charge}{电荷}守恒）？求
二者间的作用力，是吸引还是排斥？
\end{exercise}

\begin{exercise}[$\star\star$]\label{exo:g11:fundamental-interactions:9}
两个相同\omterm{def:g11:fundamental-interactions:charge}{电荷} $q = \qty{1.0e-6}{C}$ 相斥，斥力为 \qty{0.90}{N}：
它们相距多远？力降至 \qty{0.10}{N} 的位置在哪里？
\end{exercise}

\begin{exercise}[$\star\star$]\label{exo:g11:fundamental-interactions:10}
相邻\omterm{def:g11:fundamental-interactions:ladder}{分子}中的两个质子相距 $d = \qty{1.0e-10}{m}$。
求二者间的\omterm{def:g11:fundamental-interactions:charge}{库仑}斥力。在此距离下\omterm{def:g11:fundamental-interactions:strong}{强相互作用}
贡献多少？哪种相互作用使\omterm{def:g11:fundamental-interactions:ladder}{分子}结合在一起？
\end{exercise}

\begin{exercise}[$\star\star$]\label{exo:g11:fundamental-interactions:11}
两枚硬币各约含 $\num{e23}$ 个电子。将其中
\emph{百万}分之一的电子转移到另一枚上，并使两枚硬币相距 \qty{1.0}{m}。
求每枚硬币所带的\omterm{def:g11:fundamental-interactions:charge}{电荷}，再求二者间的作用力。多大质量的物体
所受\omterm{def:g10:universal-gravitation:weight}{重力}与之相等？由此说明日常物质为何呈电中性。
\end{exercise}

\begin{exercise}[$\star\star\star$]\label{exo:g11:fundamental-interactions:12}
一个铀\omterm{def:g11:fundamental-interactions:ladder}{原子核}直径为 \qty{1.4e-14}{m}，含有 $92$ 个质子。
求位于两端的两个质子之间的斥力。能否由
\omterm{def:g11:fundamental-interactions:strong}{强相互作用}束缚这一对质子？为何不能？共有多少对
相互排斥的质子对？为何重\omterm{def:g11:fundamental-interactions:ladder}{原子核}需要多余的中子，以及
为何超过某一尺寸后连中子也无法挽救它们？
\end{exercise}

\begin{exercise}[$\star\star\star$]\label{exo:g11:fundamental-interactions:13}
两个质量为 \qty{1.0}{g} 的小球，各用 \qty{50}{cm} 长的细线
系于同一点。二者带有相同\omterm{def:g11:fundamental-interactions:charge}{电荷} $q$，静止后相距 \qty{6.0}{cm}
。每个小球在\omterm{def:g10:universal-gravitation:weight}{重力}、细线\omterm{def:g10:inertia:inventory}{张力}与
电场力作用下处于平衡：证明 $F = mg\tan\theta$（$\theta$
为细线与竖直方向的夹角），并证明此处 $\tan\theta \approx 0.060$。
依次计算 $F$、$q$，再求移动的\omterm{def:g11:fundamental-interactions:elementary}{元电荷}个数。
\end{exercise}

\begin{exercise}[$\star\star\star$]\label{exo:g11:fundamental-interactions:14}
已知：$M_E = \qty{5.97e24}{kg}$，$M_M = \qty{7.35e22}{kg}$，地月
距离 $d = \qty{3.84e8}{m}$。计算地月之间的\omterm{def:g11:fundamental-interactions:four}{万有引力}
关闭引力后，若地球带电 $+Q$、月球带电 $-Q$，
要保持相同的吸引力，则 $Q$ 为多少？它相当于多少个电子？
这些电子的质量为多少（$m_e = \qty{9.11e-31}{kg}$）？请加以评论。
\end{exercise}

\begin{exercise}[$\star\star\star$]\label{exo:g11:fundamental-interactions:15}
在食盐晶体中，Na$^{+}$ 与 Cl$^{-}$ 离子（\omterm{def:g11:fundamental-interactions:charge}{电荷}为 $\pm e$）相距
$d = \qty{2.8e-10}{m}$，一个 Cl$^{-}$ 离子的质量为 \qty{5.9e-26}{kg}。
求相邻离子之间的作用力，以及
一个 Cl$^{-}$ 离子的\omterm{def:g10:universal-gravitation:weight}{重力}，并求二者之比。然后说明引力在
大尺度上为何占上风——自然界中没有超过 \qty{e4}{m} 的山峰——尽管每个化学键
的作用都比引力大十五个\omterm{def:g10:orders-of-magnitude:oom}{数量级}。
\end{exercise}

\section{周末问题：四种基本相互作用的核查}

\begin{problem}[{周末问题——为何世界既不坍缩
也不四散飞开：逐层核查四种\omterm{def:g11:fundamental-interactions:four}{基本相互作用}，直至
物质的电中性精确到 $10^{18}$ 分之二}]
\label{pb:g11:fundamental-interactions:1}
\omterm{def:g11:fundamental-interactions:ladder}{原子}不会内爆，\omterm{def:g11:fundamental-interactions:ladder}{原子核}大多能够保持稳定，月球既不会
撞向地球，也不会逃逸。本题要核查各种相互作用各自的贡献。数据：
$e = \qty{1.60e-19}{C}$,
$k = \qty{8.99e9}{N.m^2/C^2}$, $G = \qty{6.67e-11}{N.m^2/kg^2}$,
$m_p = \qty{1.67e-27}{kg}$, $m_e = \qty{9.11e-31}{kg}$,
$M_E = \qty{5.97e24}{kg}$，$M_M = \qty{7.35e22}{kg}$，地月
距离 \qty{3.84e8}{m}。

\medskip
\noindent\textbf{第一部分——阶梯。}
\begin{enumerate}
  \item 列出从\omterm{def:g11:fundamental-interactions:ladder}{夸克}到星系的各层阶梯，相邻两层
        的尺度相差一个\omterm{def:g10:orders-of-magnitude:oom}{数量级}。
  \item 缩放模型：若\omterm{def:g11:fundamental-interactions:ladder}{原子核}成为直径为 \qty{1}{cm} 的弹珠，
        \omterm{def:g11:fundamental-interactions:ladder}{原子}有多宽？
  \item \omterm{def:g11:fundamental-interactions:ladder}{原子核}占\omterm{def:g11:fundamental-interactions:ladder}{原子}体积的比例为多少？
  \item 一根头发的厚度为 \qty{70}{\micro m}，需要多少个\omterm{def:g11:fundamental-interactions:ladder}{原子}才能横跨它？
  \item 从\omterm{def:g11:fundamental-interactions:ladder}{夸克}上限（\qty{e-18}{m}）
        到星系（\qty{e21}{m}）之间相差多少个十的幂次？
\end{enumerate}

\medskip
\noindent\textbf{第二部分——电学阶梯。}
\begin{enumerate}[resume]
  \item 在氢\omterm{def:g11:fundamental-interactions:ladder}{原子}中，电子与质子之间的距离为 $d = \qty{5.3e-11}{m}$，
        求二者之间的电场力。
  \item 求二者之间的\omterm{def:g11:fundamental-interactions:four}{万有引力}，并求两种力之比。
        哪一种力将\omterm{def:g11:fundamental-interactions:ladder}{原子}束缚在一起？
  \item 若在\omterm{def:g11:fundamental-interactions:ladder}{原子}内部关闭引力，会发生什么变化？若关闭
        电的作用，又会怎样？
  \item 为什么日常生活中的“接触”力——地面对脚的支持力、\omterm{def:g10:inertia:inventory}{摩擦力}
——本质上都是伪装的电磁作用？
  \item 两枚硬币各约含 $\num{e23}$ 个电子。在二者之间转移
        每\emph{十亿}个电子中的一个，并将二者
        相隔 \qty{1.0}{m}：求\omterm{def:g11:fundamental-interactions:charge}{电荷}与作用力。结论如何？
  \item 从 \qty{e-10}{m} 到
        \qty{e4}{m} 的各层阶梯由哪种相互作用主导？为何\omterm{def:g11:fundamental-interactions:strong}{强相互作用}不参与竞争？
\end{enumerate}

\medskip
\noindent\textbf{第三部分——核阶梯。}
\begin{enumerate}[resume]
  \item \omterm{def:g11:fundamental-interactions:ladder}{原子核}中两个质子相距 \qty{1.0e-15}{m}，计算
        二者之间的电斥力。
  \item 若无其他约束，此力会使质子产生多大的加速度？
        与 $g$ 相比较。
  \item 列出\omterm{def:g11:fundamental-interactions:strong}{强相互作用}为解释以下现象所需具备的三条性质：
        为何\omterm{def:g11:fundamental-interactions:ladder}{原子核}能够存在，为何中子也被束缚，以及为何
        \omterm{def:g11:fundamental-interactions:ladder}{原子核}如此之小。
  \item 在铀\omterm{def:g11:fundamental-interactions:ladder}{原子核}（直径 \qty{1.4e-14}{m}，含 $92$ 个质子）中，
        计算位于两端的两个质子之间的斥力。能否
        由\omterm{def:g11:fundamental-interactions:strong}{强相互作用}直接束缚这一对质子？请说明
        \omterm{def:g11:fundamental-interactions:ladder}{原子核}增大时斥力为何占优，以及中子的作用。
  \item 按前文承诺，用一句话说明\omterm{def:g11:fundamental-interactions:weak}{弱相互作用}做什么。
\end{enumerate}

\medskip
\noindent\textbf{第四部分——天文阶梯与结论。}
\begin{enumerate}[resume]
  \item 计算地球与月球之间的\omterm{def:g11:fundamental-interactions:four}{万有引力}。
  \item 地球约含 $\num{1.8e51}$ 个质子（电子数相同），月球约含
        $\num{2.2e49}$ 个质子。求每个天体中未被抵消的
        质子\omterm{def:g11:fundamental-interactions:charge}{电荷}所占的比例 $f$ 为多少，
        才能使电场力与\omterm{def:g11:fundamental-interactions:four}{万有引力}相等？
  \item 第 7 题中的“弱者”\omterm{def:g11:fundamental-interactions:four}{万有引力}为何仍主宰天文学？
        请给出两个原因。
  \item 结论——每层用一句话概括（\omterm{def:g11:fundamental-interactions:ladder}{原子核}、\omterm{def:g11:fundamental-interactions:ladder}{原子}至山峰，
        行星及更大尺度）：是什么将其束缚在一起？然后回答标题中的问题：
        为何世界既不坍缩，也不四散飞开？
\end{enumerate}
\end{problem}
